% Original analytic teaching diagram, 2026-11-10.
\definecolor{VB}{HTML}{7188A3}\definecolor{VBF}{HTML}{E5EAF0}
\definecolor{VG}{HTML}{829C8B}\definecolor{VGF}{HTML}{E8EEE9}
\definecolor{VY}{HTML}{B4A16F}\definecolor{VYF}{HTML}{F1ECDD}
\definecolor{VR}{HTML}{A7787D}\definecolor{VRF}{HTML}{F0E5E6}
\definecolor{VP}{HTML}{9588A5}\definecolor{VPF}{HTML}{ECE8F0}
\definecolor{VI}{HTML}{3B4252}\definecolor{VM}{HTML}{6B7280}\definecolor{VL}{HTML}{4C4D4F}
\begin{tikzpicture}[x=1mm,y=1mm,>=Stealth,font=\figurefont\fontsize{8}{10}\selectfont,
lab/.style={text=VM,font=\figurefont\fontsize{7.5}{9}\selectfont},
bx/.style={draw=VB,fill=VBF,rounded corners=1.5pt,line width=.75pt,align=center,inner sep=3pt,minimum height=11mm,text width=25mm,text=VI},
cond/.style={lab,align=center,text=VM},
ar/.style={VL,line width=.75pt,->}]
\node[bx] (peano) at (14,48) {局部存在\\[-1pt]（可能多解）};
\node[bx] (picard) at (14,26) {局部存在且唯一\\[-1pt]Picard迭代};
\node[bx] (max) at (56,26) {最大解\\[-1pt]有限端点离开紧集};
\node[bx] (glob) at (97,26) {解延续到\\[-1pt]整个时间区间};
\draw[ar] (peano) -- (picard) node[midway,right,cond,text width=30mm]{状态方向局部Lipschitz，\\[-1pt]对时间局部一致};
\draw[ar] (picard) -- (max) node[midway,above,cond]{相容拼接};
\draw[ar] (max) -- (glob) node[midway,above,cond,text width=14mm]{至多\\[-1pt]线性增长};
\node[cond] at (14,57) {$\bm{f}$连续，有限维};
\node[bx,fill=VYF,draw=VY] (car) at (56,5) {时间可测：改用\\[-1pt]绝对连续解};
\node[bx,fill=VGF,draw=VG] (man) at (97,5) {流形：在坐标中\\[-1pt]用同一套结论};
\draw[ar,dash pattern=on 2.5pt off 1.5pt] (max.south) -- (car.north);
\draw[ar,dash pattern=on 2.5pt off 1.5pt] (max.south east) to[out=-50,in=120] (man.north);
\end{tikzpicture}
